\subsection{Coxeter 群的几何表示的不可约性}

我们沿用前面各号的记号，并假设 S 有限。

命题 7. 假设 (W,S) 是不可约的（第 IV 章，§1，第 9 号）。设 \(\mathbf{E}^0\) 是 E 的关于 \(\mathbf{B}_{\mathbf{M}}\) 与 E 正交的子空间。群 W 平凡地作用于 \(\mathbf{E}^0\) ，并且 E 的每个异于 E 且在 W 下稳定的子空间都含于 \(\mathbf{E}^0\) 。

若 \(x \in \mathrm{E}^0\) ，则 \(\sigma_s(x) = x - 2\mathrm{B}_{\mathrm{M}}(e_s, x)e_s = x\) 对所有 \(s \in \mathrm{S}\) 成立。由于 \(\mathrm{W}\) 由 \(\mathrm{S}\) 生成，由此可知 \(\mathrm{W}\) 平凡地作用于 \(\mathrm{E}^0\) 。

设 \(E'\) 是 \(E\) 的一个在 \(W\) 下稳定的子空间。设 \(s, s' \in S\) 是在图 \(\Gamma\) （即 \((W, S)\) 的图；第 IV 章，§1，第 9 号）中相连的两个元素；回忆这意味着 \(m(s, s') \geqslant 3\) 。假设 \(e_s \in E'\) 。则 \(\sigma_{s'}(e_s) \in E'\) ，且由于 \(e_{s'}\) 在 \(\sigma_{s'}(e_s)\) 中的系数非零，我们有 \(e_{s'} \in E'\) 。由于 \(\Gamma\) 连通，由此可知，若 \(E'\) 含有某个 \(e_s\) ，则它含有它们全体并与 \(E\) 重合。除这一情形外，由 §2 第 2 号命题 3 可知，对所有 \(s \in S\) ，\(E'\) 含于超平面 \(H_s\) （正交于 \(e_s\) ）。由于诸 \(H_s\) 的交等于 \(E^0\) ，命题得证。

推论. 假设 (W, S) 是不可约的。则：

\begin{enumerate}
\def\labelenumi{\alph{enumi})}
\item
  若 \(\mathrm{B_M}\) 非退化，则 W-模 E 是绝对单的。
\item
  若 \(\mathrm{B_M}\) 退化，则 W-模 E 不是半单的。
\end{enumerate}

在情形 \(a\) ，命题 7 表明 \(E\) 是单的，从而也是绝对单的（§ 2，第 1 号，命题 1）。

在情形 \(b\) ，有 \(\mathrm{E}^0 \neq 0\) ，\(\mathrm{E} \neq \mathrm{E}^0\) （因为 \(\mathrm{B}_{\mathrm{M}} \neq 0\) ），且命题 7 表明 \(\mathrm{E}^0\) 没有在 W 下稳定的补空间；因此 W-模 E 不是半单的。

